\documentclass[UTF8,11pt]{ctexart}
\usepackage[a4paper,margin=25mm]{geometry}
\usepackage{amsmath,amssymb,amsthm,mathtools}
\usepackage[all]{xy}
\usepackage{xurl}
\usepackage{hyperref}
\usepackage{enumitem}
\hypersetup{hidelinks,breaklinks=true}
\setlength{\emergencystretch}{3em}
\newtheorem*{theorem}{Theorem}
\newtheorem*{lemma}{Lemma}
\newtheorem*{proposition}{Proposition}
\newtheorem*{corollary}{Corollary}
\newtheorem*{definition}{Definition}
\newcommand{\Spec}{\operatorname{Spec}}
\newcommand{\Hom}{\operatorname{Hom}}
\newcommand{\End}{\operatorname{End}}
\newcommand{\Gal}{\operatorname{Gal}}
\newcommand{\Cl}{\operatorname{Cl}}
\newcommand{\vol}{\operatorname{vol}}
\newcommand{\covol}{\operatorname{covol}}
\newcommand{\tr}{\operatorname{tr}}
\newcommand{\im}{\operatorname{im}}
\newcommand{\id}{\operatorname{id}}
\newcommand{\coker}{\operatorname{coker}}
\newcommand{\stacksref}[2]{\href{https://stacks.math.columbia.edu/tag/#1}{#2 [#1]}}
\begin{document}
\section*{009\quad 平面直线的高重数交点上界}
\subsection*{来源与整理范围}
集体来源：MIT OpenCourseWare，MIT 18.S997 \emph{The Polynomial Method}，Fall 2012；授课教师 Lawrence D. Guth。课程注明讲义由 MIT 学生提供并获准使用，所取 PDF 未单列执笔人。Lecture 7, \emph{Crossing Numbers and the Szemer\'edi--Trotter Theorem}，Theorem 2.1，p.~3；核心支撑为 Proposition 1.2、Theorem 1.3 及后续回补的重画边论证，pp.~1--3，均在本条收录。

原讲义说明此路线遵循 L\'aszl\'o A. Sz\'ekely, ``Crossing numbers and hard Erd\H{o}s problems in discrete geometry'', \emph{Combinatorics, Probability and Computing} 6 (1997), no.~3, 353--358。此处收录的是课程讲义的证明，不冒充该论文原文。获取日期：2026-09-28。
\par\url{https://ocw.mit.edu/courses/18-s997-the-polynomial-method-fall-2012/a97e1d93e47a971e0c65113109cde508_MIT18_S997F12_lec7.pdf}
\par\url{https://ocw.mit.edu/courses/18-s997-the-polynomial-method-fall-2012/pages/lecture-notes/}

\textbf{编辑补明的使用范围（非作者原句）：}本条的图为有限简单图，直线集合中的直线互异，$k$ 为至少 $2$ 的整数。交叉数取平面中一般位置画法的最少边交叉数。原 Theorem 2.1 没有在题句中单列 $k\geq2$；这里把使用范围明示，而不把 $k=1$ 纳入本条命题。

\subsection*{作者原命题与人类证明}
\paragraph{Crossing number estimates}
\begin{proposition}[1.2]
The crossing number of $G$ is at least $E-3V$.
\end{proposition}
\begin{proof}
Let $k(G)$ be the crossing number of $G$. Embed $G$ in the plane with $k(G)$ crossings. By removing at most $k(G)$ edges, we get a planar graph $G'$ with $E'=E-k$ edges and $V'\leq V$ vertices. We see $0\geq E'-3V'\geq E-k-3V$.
\end{proof}

\begin{theorem}[1.3]
If $G$ is a graph with $E$ edges and $V$ vertices, and $E\geq4V$, then the crossing number of $G$ is at least $(1/64)E^3V^{-2}$.
\end{theorem}
This theorem was proven by several authors before Sz\'ekely, but we give his proof.
\begin{proof}
Let $p$ be a number between $0$ and $1$ which we choose below. Let $G'$ be a random subgraph of $G$ formed by including each vertex of $G$ independently with probability $p$. We include an edge of $G$ in $G'$ if its endpoints are in $G'$.

We consider the expected values for the number of vertices and edges in $G'$. The expected value of $V'$ is $pV$. The expected value of $E'$ is $p^2E$. For every subgraph $G'\subset G$, the crossing number of $G'$ is at least $E'-3V'$. Therefore, the expected value of the crossing number of $G'$ is at least $p^2E-3pV$.

On the other hand, we give an upper bound on the expected crossing number of $G'$ as follows. Let $k=k(G)$ be the crossing number of $G$. Let $F:G\to\mathbb R^2$ be a legal embedding with $k$ crossings. We claim that each crossing of $F$ involves two disjoint edges. In other words, two edges that share a vertex don't cross. We come back to the claim at the end. By restricting $F$ to $G'$, we get an embedding of $G'$ with $p^4k$ crossings on average. This is because each crossing involves four vertices, and it appears as a crossing of $F(G')$ only if all four vertices are included in $G'$. (If $F$ had a crossing involving two edges containing a common vertex, then it would appear with the much higher probability $p^3$.) Therefore, the expected value of the crossing number of $G'$ is at most $p^4k$.

Comparing our upper and lower bounds, we see that $p^4k\geq p^2E-3pV$, and so we get the following lower bound for $k$.
\[
k\geq p^{-2}E-3p^{-3}V.
\]
We can now choose $p$ to optimize the right-hand side. We choose $p=4V/E$, and we have $p\leq1$ since we assumed $4V\leq E$. Plugging in we get $k\geq(1/64)E^3V^{-2}$.

To finish the proof, we just have to check the claim that $F$ has no crossings of edges that share a vertex. Given any map with such a crossing, we explain how to modify it to reduce the crossing number. Say that $F(e_1)$ and $F(e_2)$ each leave $F(v)$ and cross at $x$. (If they cross several times, then let $x$ be the last crossing.) We modify $F$ as follows. Suppose that $F(e_1)$ crosses $k_1$ other edges on the way from $F(v)$ to $x$ and that $F(e_2)$ crosses $k_2$ other edges on the way from $F(v)$ to $x$. We choose the labelling so that $k_1\leq k_2$. Then we modify $F$ on the edge $e_2$, making $F(e_2)$ follow parallel to $F(e_1)$ until $x$ and then rejoin its original course at $x$, so that $F(e_1)$ and $F(e_2)$ never cross. This operation reduces the crossing number of $x$, and so a minimal map $F$ has no such crossings.
\end{proof}

\paragraph{The Szemer\'edi--Trotter theorem}
\begin{theorem}[2.1]
Let $\mathcal L$ be a set of $L$ lines in the plane. Let $P_k$ be the set of points that lie on at least $k$ lines of $\mathcal L$. Then the number of points in $P_k$ is at most
\[
\max(2Lk^{-1},2^9L^2k^{-3}).
\]
\end{theorem}
\begin{proof}
Using the lines and points, we make a graph mapped into the plane. The vertices of our graph $G$ are the points of $P_k$. We join two vertices with an edge of $G$ if the two points are two consecutive points of $P_k$ on a line $l\in\mathcal L$. This graph is not embedded, but the crossing number of our map is at most $\binom L2\leq L^2$, since each crossing of the graph $G$ must correspond to an intersection of two lines of $\mathcal L$.

We will count the vertices and edges of the graph $G$ and apply the crossing number theorem. The number of vertices of our graph is $V=|P_k|$. The number of edges of our graph is $kV-L$. (At first sight, each vertex should be adjacent to $2k$ edges which would give $kV$ edges. But on each line $l\in\mathcal L$, the first and last vertices are adjacent to one less edge than this initial count.) As long as $E\geq4V$, we can apply the crossing number theorem and it gives
\[
L^2\geq(1/64)(kV-L)^3V^{-2}.
\]
Either $V\leq2L/k$, or else $kV-L\geq(1/2)kV$. In the former case, we are done. In the latter case, we have $L^2\geq2^{-9}k^3V$, which means $V\leq2^9L^2k^{-3}$.

On the other hand, if $E<4V$, we have $kV-L\leq4V$, and hence $V\leq\frac{L}{k-4}$. As long as $k\geq8$, this implies $V\leq2L/k$, and we are done. Finally, for $k<8$, the trivial bound
\[
|P_k|\leq\binom L2/\binom k2\leq2L^2k^{-2}\leq2^9L^2k^{-3}.
\]
\end{proof}


\subsection*{整理说明（非作者正文）}
本条主命题为 Theorem 2.1；Theorem 1.3 及 Proposition 1.2 作为其核心证明链合并收录，不各计一道题。允许的标准先修是有限简单平面图的 Euler 公式推论 $E-3V\leq0$（原 Proposition 1.1），以及有限概率空间中的期望线性性。

保留作者的先证明交叉数界、再返回直线交点计数的顺序。回补的 ``two edges that share a vertex don't cross'' 论证全文已收录，没有停在 ``We come back to the claim''。略去原文的动机、完全图例子和不承担证明的历史文字。

\textbf{原文措辞保留记录：}原 Theorem 2.1 的证明把边数写为 ``is $kV-L$''，本转录照录，不将其静默改写成等式之外的另一版本；原文的 ``crossing number of $x$'' 也照录。这些保留不表示对其措辞作了数学认证。直线集合记号排作 $\mathcal L$，与数量 $L$ 区分。已取得 PDF 文本，但页面截图接口失败，未声称逐页图像核对，未保留下载 PDF 原件。

\subsection*{可修改的简短标注（非作者正文）}
$c$：平面直线配置的高重数交点计数。

$a$：由共线相邻点构造图；平面图 Euler 不等式；交叉数；独立顶点抽样；期望计数；局部重画边；参数优化。

\texttt{cross\_theory\_status = yes}。几何关联被编码成平面中画出的图，拓扑交叉约束经随机诱导子图放大为三次边数界，再译回高重数交点上界。位置：Theorem 1.3 的随机抽样证明及 Theorem 2.1 的图构造，pp.~2--3。标注为非穷尽、可修改初稿。
\subsection*{许可与修改记录}
材料来自 MIT OpenCourseWare；原作者与课程见上文。按该站所示 \href{https://creativecommons.org/licenses/by-nc-sa/4.0/}{CC BY-NC-SA 4.0} 许可复用。本转录及其附加整理采用同一许可；不得据此声称作者或 MIT 认可本次整理。2026-09-28：ChatGPT 转录为 TeX，加入题号、来源、整理说明和初稿标注；数学正文保持作者语言及顺序。

\end{document}
